% ============================================================================
% CHAPTER 2: SYSTEM ARCHITECTURE
% ============================================================================

\chapter{System Architecture}
\label{ch:architecture}

This chapter provides a detailed description of the system architecture, including the control pipeline, data flow between modules, coordinate systems, and units of measurement.

\section{Control Pipeline Overview}

The system implements a hierarchical control architecture where each layer operates at a different level of abstraction:

\begin{figure}[htbp]
    \centering
    \begin{tikzpicture}[
        scale=0.85,
        transform shape,
        layer/.style={
            draw,
            rounded corners=5pt,
            minimum width=14cm,
            minimum height=1.8cm,
            align=center
        }
    ]
        % Layers
        \node[layer, fill=contractpurple!20] (planning) at (0,6) {
            \textbf{Planning Layer}\\
            \small Horizon Planner $\bullet$ Mission Management $\bullet$ Contract Verification
        };

        \node[layer, fill=supervisororange!20] (supervision) at (0,4) {
            \textbf{Supervision Layer}\\
            \small Flight Mode Management $\bullet$ Controller Selection $\bullet$ Setpoint Generation
        };

        \node[layer, fill=pidblue!20] (control) at (0,2) {
            \textbf{Control Layer}\\
            \small PID Controller $\bullet$ H-$\infty$ Controller $\bullet$ Attitude/Rate Control
        };

        \node[layer, fill=cbforange!20] (safety) at (0,0) {
            \textbf{Safety Layer}\\
            \small CBF Filter $\bullet$ Hard Constraints $\bullet$ Emergency Intervention
        };

        \node[layer, fill=gray!20] (physical) at (0,-2) {
            \textbf{Physical Layer}\\
            \small Actuators $\bullet$ Sensors $\bullet$ Dynamics
        };

        % Arrows between layers
        \draw[-{Stealth}, thick] (planning.south) -- (supervision.north)
            node[midway, right, font=\small] {controller recommendation};
        \draw[-{Stealth}, thick] (supervision.south) -- (control.north)
            node[midway, right, font=\small] {setpoints};
        \draw[-{Stealth}, thick] (control.south) -- (safety.north)
            node[midway, right, font=\small] {nominal control};
        \draw[-{Stealth}, thick] (safety.south) -- (physical.north)
            node[midway, right, font=\small] {safe control};

        % Feedback
        \draw[-{Stealth}, thick, dashed]
            (physical.west) -- ++(-1,0) |- (planning.west)
            node[pos=0.25, left, font=\small] {state feedback};

        % Time scales
        \node[right=1cm of planning, font=\footnotesize] {$\sim$0.5s update};
        \node[right=1cm of supervision, font=\footnotesize] {$\sim$0.1s update};
        \node[right=1cm of control, font=\footnotesize] {50 Hz (20ms)};
        \node[right=1cm of safety, font=\footnotesize] {50 Hz (20ms)};
    \end{tikzpicture}
    \caption{Hierarchical control architecture with different time scales}
    \label{fig:layers}
\end{figure}

\section{Coordinate System: NED Frame}

The system uses the \textbf{North-East-Down (NED)} coordinate frame, which is the standard for aerospace applications.

\begin{figure}[htbp]
    \centering
    \begin{tikzpicture}[scale=1.2]
        % Ground plane
        \fill[brown!20] (-3,-0.5) -- (3,-0.5) -- (4,1) -- (-2,1) -- cycle;
        \draw[brown!50, thick] (-3,-0.5) -- (3,-0.5) -- (4,1) -- (-2,1) -- cycle;

        % Drone body
        \node[circle, fill=gray!50, minimum size=0.8cm] (drone) at (0,2) {};

        % Coordinate axes
        \draw[-{Stealth}, very thick, red] (drone.center) -- ++(2,0)
            node[right, font=\bfseries] {X (North)};
        \draw[-{Stealth}, very thick, green!60!black] (drone.center) -- ++(1.5,-1)
            node[below right, font=\bfseries] {Y (East)};
        \draw[-{Stealth}, very thick, blue] (drone.center) -- ++(0,-1.5)
            node[below, font=\bfseries] {Z (Down)};

        % Altitude annotations
        \draw[<->, thick] (-2.5,-0.5) -- (-2.5,2)
            node[midway, left, font=\small] {altitude};

        % Ground level
        \node[font=\footnotesize] at (3.5,0.2) {Ground ($z=0$)};

        % Key points
        \node[font=\footnotesize, text=blue] at (1.5,0.5) {$z > 0$: below ground};
        \node[font=\footnotesize, text=blue] at (1.5,3) {$z < 0$: above ground};
    \end{tikzpicture}
    \caption{NED coordinate frame definition}
    \label{fig:ned}
\end{figure}

\warningbox{
In the NED frame, \textbf{negative} $z$ values represent altitude above ground. For example:
\begin{itemize}
    \item $z = -5.0$ m means 5 meters \textbf{above} ground
    \item $z = 0$ means at ground level
    \item $z > 0$ means \textbf{below} ground (collision!)
\end{itemize}
This convention is critical for correct altitude control and safety constraints.
}

\section{Complete State Vector}

The system maintains a 12-dimensional state vector:

\begin{equation}
    \mathbf{x} = \begin{bmatrix}
        x \\ y \\ z \\ v_x \\ v_y \\ v_z \\ \phi \\ \theta \\ \psi \\ p \\ q \\ r
    \end{bmatrix} \in \mathbb{R}^{12}
\end{equation}

\begin{table}[htbp]
    \centering
    \caption{State vector components with units and typical ranges}
    \label{tab:state_vector}
    \begin{tabular}{@{}clllc@{}}
        \toprule
        \textbf{Index} & \textbf{Symbol} & \textbf{Description} & \textbf{Unit} & \textbf{Range} \\
        \midrule
        1--3 & $x, y, z$ & Position (NED) & \si{\meter} & $\pm 100$ \\
        4--6 & $v_x, v_y, v_z$ & Velocity (NED) & \si{\meter\per\second} & $\pm 15$ \\
        7 & $\phi$ & Roll angle & \si{\radian} & $\pm\pi$ \\
        8 & $\theta$ & Pitch angle & \si{\radian} & $\pm\pi/2$ \\
        9 & $\psi$ & Yaw angle & \si{\radian} & $\pm\pi$ \\
        10--12 & $p, q, r$ & Angular rates (body) & \si{\radian\per\second} & $\pm 3$ \\
        \bottomrule
    \end{tabular}
\end{table}

\section{Data Flow Specification}

The following table specifies the exact data flowing between each module:

\begin{table}[htbp]
    \centering
    \caption{Inter-module data flow specification}
    \label{tab:dataflow}
    \small
    \begin{tabularx}{\textwidth}{@{}llXl@{}}
        \toprule
        \textbf{Source} & \textbf{Destination} & \textbf{Data} & \textbf{Rate} \\
        \midrule
        GPS Sensor & EKF & Position $[x,y,z]$, Velocity $[v_x,v_y,v_z]$, satellites, HDOP & 10 Hz \\
        IMU Sensor & EKF & Attitude $[\phi,\theta,\psi]$, Angular rates $[p,q,r]$ & 200 Hz \\
        Wind Sensor & Planner & Wind speed, direction & 10 Hz \\
        Battery & Monitor & Voltage, current & 1 Hz \\
        \midrule
        EKF & Planner & State estimate $\hat{\mathbf{x}}$, covariance $\mathbf{P}$ & 50 Hz \\
        EKF & Supervisor & State estimate $\hat{\mathbf{x}}$, GPS validity & 50 Hz \\
        \midrule
        Planner & Supervisor & Recommended controller, safety margin, plan status & 2 Hz \\
        Supervisor & Controller & Setpoint (position, velocity, yaw) & 50 Hz \\
        Supervisor & Controller & Controller selection (PID/H-$\infty$) & As needed \\
        \midrule
        Controller & CBF & Thrust $T$, torques $[\tau_\phi, \tau_\theta, \tau_\psi]$ & 50 Hz \\
        Controller & CBF & Desired attitude $[\phi_d, \theta_d, \psi_d]$ & 50 Hz \\
        \midrule
        CBF & Actuators & Safe thrust, safe torques & 50 Hz \\
        \bottomrule
    \end{tabularx}
\end{table}

\section{Module Interfaces}

Each module exposes a well-defined interface with typed inputs and outputs:

\subsection{Sensor Interface}

\begin{lstlisting}[caption={Sensor data structure}]
sensors = {
    'gps': {
        'position': np.array([x, y, z]),      # meters (NED)
        'velocity': np.array([vx, vy, vz]),   # m/s (NED)
        'valid': bool,                         # True if GPS lock
        'satellites': int,                     # Number of satellites
        'hdop': float                          # Horizontal DOP
    },
    'imu': {
        'attitude': np.array([phi, theta, psi]),  # radians
        'rates': np.array([p, q, r]),             # rad/s
        'valid': bool,
        'calibrated': bool,
        'temperature': float                       # Celsius
    },
    'wind': {
        'speed': float,     # m/s
        'direction': float  # radians
    },
    'battery': {
        'voltage': float    # Volts
    }
}
\end{lstlisting}

\subsection{Control Interface}

\begin{lstlisting}[caption={Control command structure}]
control = {
    'thrust': float,                    # Normalized [0, 1]
    'torques': np.array([tx, ty, tz]),  # Nm
    'desired_attitude': np.array([...]), # radians
    'desired_rates': np.array([...])     # rad/s
}
\end{lstlisting}

\subsection{Setpoint Interface}

\begin{lstlisting}[caption={Setpoint structure from supervisor}]
setpoint = {
    'position': np.array([x, y, z]),    # meters (NED)
    'velocity': np.array([vx, vy, vz]), # m/s
    'yaw': float                         # radians
}
\end{lstlisting}

\section{Timing and Synchronization}

The system operates with the following timing constraints:

\begin{figure}[htbp]
    \centering
    \begin{tikzpicture}[scale=0.8]
        % Timeline
        \draw[-{Stealth}, thick] (0,0) -- (14,0) node[right] {Time (ms)};

        % Tick marks
        \foreach \x in {0,2,4,6,8,10,12} {
            \draw (\x,0.1) -- (\x,-0.1) node[below, font=\footnotesize] {\x0};
        }

        % Control loop
        \foreach \x in {0,2,4,6,8,10,12} {
            \fill[pidblue] (\x,1) rectangle (\x+0.3,1.5);
        }
        \node[right, font=\small] at (13,1.25) {Control (20ms)};

        % Planner updates
        \foreach \x in {0,10} {
            \fill[contractpurple] (\x,2) rectangle (\x+1,2.5);
        }
        \node[right, font=\small] at (13,2.25) {Planner (500ms)};

        % Sensor reads
        \foreach \x in {0,1,2,3,4,5,6,7,8,9,10,11,12} {
            \fill[sensorblue] (\x,3) rectangle (\x+0.15,3.5);
        }
        \node[right, font=\small] at (13,3.25) {IMU (100Hz)};

        \foreach \x in {0,2,4,6,8,10,12} {
            \fill[ekfgreen] (\x,4) rectangle (\x+0.2,4.5);
        }
        \node[right, font=\small] at (13,4.25) {GPS (50Hz)};
    \end{tikzpicture}
    \caption{Timing diagram showing multi-rate operation}
    \label{fig:timing}
\end{figure}

\section{File Organization}

The implementation follows a modular file structure:

\begin{lstlisting}[language=bash, caption={Project directory structure}]
contract-based-uav-control/
|-- src/
|   |-- adaptive_control_system.py     # Main integration
|   |
|   |-- contracts/
|   |   |-- contract_framework.py      # A/G contract definitions
|   |
|   |-- estimation/
|   |   |-- contract_ekf.py            # Contract-aware EKF
|   |
|   |-- control/
|   |   |-- controllers.py             # PID, H-infinity
|   |   |-- flight_mode_supervisor.py  # Mode management
|   |
|   |-- planning/
|   |   |-- horizon_planner.py         # Pacti-based planner
|   |   |-- integrated_planner.py      # Control integration
|   |   |-- pacti_contracts.py         # Pacti interface
|   |
|   |-- safety/
|       |-- cbf_filter.py              # CBF safety filter
|
|-- demo_waypoint_mission.py           # Demo script
|-- tests/                             # Unit tests
|-- docs/                              # Documentation
\end{lstlisting}
